\documentclass[11pt]{article}
\usepackage{series}

\hypersetup{
  pdftitle={Canonical Coherent Kernels from MTT Fixed-Point Data},
  pdfauthor={Peter Nero}
}

\title{\textbf{Canonical Coherent Kernels from MTT Fixed-Point Data}\\
\large From the Delta Dictionary to Execution-Level Corrections}
\author{Peter Nero}
\date{\today}

\begin{document}
\maketitle

\begin{abstract}
Earlier papers in the delta--projection sequence established a structural dictionary:
Dirac delta distributions arise as singular limits of finite kernels, shells, filters, windows,
measurement effects, and thin layers. That dictionary is mathematically legitimate but does
not by itself determine new physics, because an arbitrary finite kernel is only a regulator.
The execution-level problem is the reverse direction: derive the finite kernel from MTT data.

This paper proves the canonical fixed-point version of that reverse direction. In an admissible
MTT fixed-point regime, the data consist of a nonnegative self-adjoint linearized operator
\(A\), a bounded coherent projector \(P=\Pi_{\mathrm{coh}}\), a complementary incoherent projector
\(Q=I-P\), a spectral gap on the incoherent sector, and a proper-time/coherence scale
\(\tau>0\). These data canonically define the coherent kernel
\[
        K_{\mathrm{MTT}}(x,y;\tau)
        =
        \langle x|\,P e^{-\tau A}P\,|y\rangle .
\]
Equivalently, in a spectral representation,
\[
        K_{\mathrm{MTT}}(x,y;\tau)
        =
        \sum_{n\in{\rm coh}} e^{-\tau\lambda_n}
        \phi_n(x)\phi_n^\ast(y).
\]
We prove that this kernel is smooth for positive proper time, is functorial under unitary
re-encoding, acts as a controlled sector-identity on the coherent sector, suppresses incoherent
components by the fixed-point spectral gap, and converges to the usual Dirac delta only in the
joint limit of complete coherent bandwidth and vanishing proper-time width.

Thus the delta replacement is not arbitrary:
\[
        \delta(x-y)
        \quad\leadsto\quad
        K_{\mathrm{MTT}}(x,y;\tau),
\]
with finite corrections controlled by \((A,P,\tau,\lambda^\ast)\). This converts the previous
delta dictionary into an execution rule: MTT fixed-point data determine the finite kernel,
and the finite kernel determines the deviation from the sharp standard object.
\end{abstract}

\seriestagline

\tableofcontents

\section{Purpose and claim discipline}

The earlier delta--projection papers established statements of the form
\[
        K_\epsilon \longrightarrow \delta
\]
or
\[
        \text{finite object} \longrightarrow \text{sharp standard object}.
\]
Those statements explain why delta distributions recur across physics, but they do not yet
select a physical finite kernel. The missing execution-level direction is
\[
        \text{MTT fixed-point data}
        \Longrightarrow
        K_{\mathrm{coh}}
        \Longrightarrow
        \text{finite correction}.
\]

This paper proves that direction in the fixed-point regime.

\begin{quote}
\textbf{Central claim.}
In an admissible fixed-point regime, the finite replacement of the Dirac delta is not chosen
freely. It is the projected heat/proper-time kernel determined by the same operator and
projector data that define coherent stabilization.
\end{quote}

\subsection*{What is proved}

We prove, in a standard compact spectral setting and in the corresponding Hilbert-space
operator language, that fixed-point data \((A,P,\tau)\) determine a canonical bounded/smoothing
operator
\[
        B_\tau:=P e^{-\tau A}P
\]
and hence a kernel \(K_{\mathrm{MTT}}\). We give explicit error estimates comparing \(B_\tau\)
with the full identity and with the coherent-sector identity.

\subsection*{What is not claimed}

We do not claim here that all numerical values of \(\tau\), spectral gaps, or carrier-induced
operators have been computed in every physical sector. The point is sharper: once a fixed-point
sector supplies \((A,P,\tau)\), the finite kernel is fixed by functional calculus and is no
longer arbitrary.

\section{Analytic setting}

Let \(X\) be a compact smooth Riemannian manifold without boundary, or a compact manifold
with boundary equipped with a self-adjoint elliptic boundary condition. More generally, one
may replace scalar functions by sections of a Hermitian vector bundle; for readability we use
the scalar notation.

Let
\[
        A\ge 0
\]
be a nonnegative self-adjoint Laplace-type operator on \(L^2(X)\) with compact resolvent.
Let
\[
        A\phi_n=\lambda_n\phi_n,\qquad
        0\le \lambda_0\le \lambda_1\le \cdots,
\]
be an orthonormal spectral resolution, with eigenvalues repeated according to multiplicity.

\begin{assumption}[Fixed-point spectral data]\label{ass:fpdata}
We are given:
\begin{enumerate}[label=(\roman*)]
\item a bounded orthogonal projector \(P\) on \(L^2(X)\);
\item \(Q:=I-P\);
\item commutation \(PA=AP\) on \(\operatorname{Dom}(A)\);
\item an incoherent-sector gap
\[
        A|_{\operatorname{Ran}Q}\ge \lambda^\ast Q
        \quad\text{for some}\quad \lambda^\ast>0;
\]
\item a proper-time/coherence scale \(\tau>0\).
\end{enumerate}
\end{assumption}

\begin{remark}
In the fixed-point papers \(P\) is the coherent projector \(\Pi_{\mathrm{coh}}\), often defined by
Riesz spectral calculus on the fiber/internal operator. The gap \(\lambda^\ast\) is the uniform
positive spectral gap above the coherent modes. The operator \(A\) is the linearized damping
or elliptic operator controlling stabilization near the fixed-point sector.
\end{remark}

\begin{definition}[Canonical MTT coherent operator]\label{def:Btau}
Under \cref{ass:fpdata}, define
\[
        B_\tau:=P e^{-\tau A}P .
\]
When \(B_\tau\) has an integral kernel, we write
\[
        K_{\mathrm{MTT}}(x,y;\tau)
        :=
        \langle x|B_\tau|y\rangle .
\]
\end{definition}

Because \(P\) commutes with \(A\), the spectral form is transparent. If \(P\phi_n=p_n\phi_n\)
with \(p_n\in\{0,1\}\), then
\[
        B_\tau f
        =
        \sum_{n}p_n e^{-\tau\lambda_n}\langle \phi_n,f\rangle\phi_n,
\]
and the kernel is
\[
        K_{\mathrm{MTT}}(x,y;\tau)
        =
        \sum_{n:p_n=1}
        e^{-\tau\lambda_n}\phi_n(x)\phi_n^\ast(y).
\]

\section{Existence, smoothness, and covariance of the canonical kernel}

\begin{theorem}[Kernel existence and smoothing]\label{thm:smoothing}
Assume \cref{ass:fpdata}. For every \(\tau>0\), \(e^{-\tau A}\) is a smoothing trace-class
operator. Consequently \(B_\tau=P e^{-\tau A}P\) is bounded and trace class. If \(P\) is a
spectral projector of \(A\), then \(B_\tau\) has a smooth kernel
\[
        K_{\mathrm{MTT}}(\cdot,\cdot;\tau)\in C^\infty(X\times X).
\]
If \(P\) is finite rank, the same conclusion holds even at \(\tau=0\), with \(B_0=P\).
\end{theorem}

\begin{proof}
Since \(A\) is a nonnegative self-adjoint elliptic operator with compact resolvent,
the spectral theorem gives
\[
        e^{-\tau A}f
        =
        \sum_n e^{-\tau\lambda_n}\langle \phi_n,f\rangle\phi_n.
\]
For every \(N\ge 0\),
\[
        A^N e^{-\tau A}
\]
is bounded, because \(\lambda^N e^{-\tau\lambda}\) is bounded on \([0,\infty)\).
Elliptic regularity then implies that \(e^{-\tau A}\) maps distributions to smooth functions.
On a compact manifold, \(e^{-\tau A}\) is trace class for \(\tau>0\), and has a smooth heat
kernel. Multiplication on both sides by the bounded spectral projector \(P\), which commutes
with \(A\), preserves trace class and smoothness of the kernel. If \(P\) is finite rank,
\(P\) is a finite sum of smooth eigenfunction projectors, hence has a smooth kernel at
\(\tau=0\). 
\end{proof}

\begin{theorem}[Re-encoding covariance]\label{thm:covariance}
Let \(U:L^2(X)\to L^2(X')\) be a unitary re-encoding. Define
\[
        A'=UAU^{-1},\qquad P'=UPU^{-1}.
\]
Then the canonical coherent operator transforms as
\[
        B'_\tau=P'e^{-\tau A'}P'=UB_\tau U^{-1}.
\]
Thus the construction of \(B_\tau\) depends only on the fixed-point operator data up to
unitary re-encoding.
\end{theorem}

\begin{proof}
Functional calculus is covariant under unitary conjugation:
\[
        e^{-\tau A'}=e^{-\tau UAU^{-1}}=Ue^{-\tau A}U^{-1}.
\]
Therefore
\[
        P'e^{-\tau A'}P'
        =
        (UPU^{-1})(Ue^{-\tau A}U^{-1})(UPU^{-1})
        =
        U(Pe^{-\tau A}P)U^{-1}.
\]
\end{proof}

\section{Sector identity and error estimates}

The canonical kernel should not be confused with the full Dirac identity kernel. It is an
identity-like object on the retained coherent sector, with explicit error terms.

\begin{assumption}[Coherent bandwidth]\label{ass:band}
For some \(\Lambda_{\mathrm{coh}}\ge 0\),
\[
        \sigma(A|_{\operatorname{Ran}P})\subset[0,\Lambda_{\mathrm{coh}}].
\]
\end{assumption}

\begin{remark}
If \(P\) is the harmonic zero-mode projector, then \(\Lambda_{\mathrm{coh}}=0\) and
\(B_\tau=P\) exactly. If \(P\) retains a finite coherent band, then
\(\Lambda_{\mathrm{coh}}\) is the largest retained coherent eigenvalue.
\end{remark}

\begin{theorem}[Coherent-sector identity estimate]\label{thm:sector}
Assume \cref{ass:fpdata,ass:band}. For every \(f\in L^2(X)\),
\[
        \norm{Pf-B_\tau f}_{L^2}
        \le
        \bigl(1-e^{-\tau\Lambda_{\mathrm{coh}}}\bigr)\norm{Pf}_{L^2}.
\]
Consequently,
\[
        \norm{f-B_\tau f}_{L^2}
        \le
        \norm{Qf}_{L^2}
        +
        \bigl(1-e^{-\tau\Lambda_{\mathrm{coh}}}\bigr)\norm{Pf}_{L^2}.
\]
\end{theorem}

\begin{proof}
Since \(P\) commutes with \(A\), \(B_\tau f=e^{-\tau A}Pf\). Thus
\[
        Pf-B_\tau f
        =
        (I-e^{-\tau A})Pf.
\]
On \(\operatorname{Ran}P\), the spectral values of \(A\) lie in
\([0,\Lambda_{\mathrm{coh}}]\). Hence
\[
        \norm{(I-e^{-\tau A})P}_{L^2\to L^2}
        =
        \sup_{\lambda\in\sigma(A|_{\operatorname{Ran}P})}
        |1-e^{-\tau\lambda}|
        \le
        1-e^{-\tau\Lambda_{\mathrm{coh}}}.
\]
This proves the first estimate. The second follows from
\[
        f-B_\tau f
        =
        Qf+(Pf-B_\tau f)
\]
and the triangle inequality.
\end{proof}

\begin{corollary}[Small-\(\tau\) sector correction]\label{cor:smalltau}
Under \cref{ass:fpdata,ass:band},
\[
        \norm{Pf-B_\tau f}_{L^2}
        \le
        \tau\Lambda_{\mathrm{coh}}\norm{Pf}_{L^2}.
\]
\end{corollary}

\begin{proof}
Use \(1-e^{-x}\le x\) for \(x\ge0\).
\end{proof}

\begin{theorem}[Incoherent damping estimate]\label{thm:gap}
Assume \cref{ass:fpdata}. Then
\[
        \norm{e^{-\tau A}Q}_{L^2\to L^2}
        \le
        e^{-\tau\lambda^\ast}.
\]
More generally, for \(s\ge0\),
\[
        \norm{A^{s/2}e^{-\tau A}Q}_{L^2\to L^2}
        \le
        \sup_{\lambda\ge\lambda^\ast}\lambda^{s/2}e^{-\tau\lambda}.
\]
\end{theorem}

\begin{proof}
On \(\operatorname{Ran}Q\), the spectral theorem and the gap assumption give
\(\lambda\ge\lambda^\ast\). Therefore
\[
        \norm{e^{-\tau A}Q}
        =
        \sup_{\lambda\in\sigma(A|_{\operatorname{Ran}Q})}e^{-\tau\lambda}
        \le e^{-\tau\lambda^\ast}.
\]
The higher estimate is identical after multiplying the spectral multiplier by
\(\lambda^{s/2}\).
\end{proof}

\section{Recovery of the Dirac delta}

The canonical MTT kernel is finite. The Dirac delta is recovered only under an idealizing
limit.

Let \(P_N:=\mathbf 1_{[0,\Lambda_N]}(A)\) be an increasing family of spectral projectors
with \(\Lambda_N\to\infty\). Let
\[
        B_{N,\tau}:=P_Ne^{-\tau A}P_N.
\]

\begin{theorem}[Joint sharp limit]\label{thm:sharp}
Let \(f\in C^\infty(X)\). Suppose \(\Lambda_N\to\infty\) and \(\tau_N\downarrow0\). Then
\[
        B_{N,\tau_N}f\to f
\]
in \(C^\infty(X)\). Equivalently, the kernels of \(B_{N,\tau_N}\) converge to
\(\delta(x-y)\) in the distributional sense on \(X\times X\).
\end{theorem}

\begin{proof}
Write
\[
        f-B_{N,\tau_N}f
        =
        (I-P_N)f
        +
        P_N(I-e^{-\tau_N A})P_Nf.
\]
The first term tends to zero in all Sobolev norms because spectral projectors converge
strongly to the identity on smooth functions, and smooth spectral coefficients decay rapidly.
For the second term, fix a Sobolev index \(s\). Choose \(M>s\). Since \(f\in C^\infty\),
\[
        \sum_n (1+\lambda_n)^M |\langle \phi_n,f\rangle|^2<\infty.
\]
For each fixed \(n\), \(1-e^{-\tau_N\lambda_n}\to0\). The summand is dominated by a
multiple of \((1+\lambda_n)^M|\langle\phi_n,f\rangle|^2\) for \(N\) large enough and
\(M\) chosen above \(s\). Dominated convergence gives convergence to zero in \(H^s\).
Since \(s\) is arbitrary, Sobolev embedding gives convergence in \(C^\infty\). The
distributional kernel statement is the usual kernel form of convergence to the identity.
\end{proof}

\begin{corollary}[Meaning of the delta limit]\label{cor:meaning}
The Dirac delta is not the fixed-point kernel itself. It is the joint idealization in which:
\begin{enumerate}[label=(\roman*)]
\item coherent bandwidth becomes complete;
\item proper-time width tends to zero;
\item the incoherent complement is no longer discarded.
\end{enumerate}
In finite-capacity MTT, the physically selected object is \(K_{\mathrm{MTT}}\), not
\(\delta\).
\end{corollary}

\section{Execution rule for downstream delta replacements}

We can now state the execution rule.

\begin{definition}[MTT delta replacement]\label{def:replacement}
In a fixed-point regime with data \((A,P,\tau)\), every downstream occurrence of an
identity/source/constraint delta that belongs to the same coherent chart is replaced by the
canonical kernel
\[
        \delta(x-y)
        \quad\rightsquigarrow\quad
        K_{\mathrm{MTT}}(x,y;\tau)
        =
        \langle x|Pe^{-\tau A}P|y\rangle.
\]
\end{definition}

This is not an arbitrary smoothing prescription. It is determined by the fixed-point data.

\begin{theorem}[Uniqueness within the heat/proper-time class]\label{thm:heatunique}
Among positive self-adjoint contraction semigroups generated by the fixed-point operator
\(A\), the one-parameter family \(e^{-\tau A}\) is uniquely determined by \(A\). Therefore,
once \(P\) and \(\tau\) are fixed, the operator \(B_\tau=Pe^{-\tau A}P\) is unique.
\end{theorem}

\begin{proof}
By the spectral theorem, a nonnegative self-adjoint operator \(A\) determines a unique
strongly continuous contraction semigroup \(T_\tau=e^{-\tau A}\). Conversely, the
generator of such a semigroup is unique by the Hille--Yosida theorem. Therefore the
proper-time filter associated with the fixed-point generator is fixed. Multiplying on the
left and right by the fixed coherent projector \(P\) gives a unique \(B_\tau\).
\end{proof}

\begin{remark}
The theorem does not say that no other regulator can be written down. It says that if the
kernel is required to be the heat/proper-time filter generated by the same operator governing
MTT stabilization, then the kernel is fixed by the fixed-point data.
\end{remark}

\section{Finite corrections in standard structures}

\subsection{Green functions}

Let \(L\) be a positive self-adjoint invertible operator commuting with \(A\) and \(P\). The
ordinary point-source Green equation
\[
        LG=I
\]
has kernel form
\[
        LG(x,y)=\delta(x-y).
\]
The MTT coherent Green operator is
\[
        G_\tau^{\mathrm{MTT}}:=L^{-1}B_\tau.
\]
It satisfies
\[
        LG_\tau^{\mathrm{MTT}}=B_\tau,
\]
or in kernels,
\[
        L_xG_\tau^{\mathrm{MTT}}(x,y)=K_{\mathrm{MTT}}(x,y;\tau).
\]

\begin{proposition}[Green-function correction]\label{prop:green}
Assume \(L^{-1}\) is bounded on the Hilbert space under consideration. Then
\[
        \norm{L^{-1}-G_\tau^{\mathrm{MTT}}}
        \le
        \norm{L^{-1}}\,
        \norm{I-B_\tau}.
\]
On a coherent-band input \(f=Pf\),
\[
        \norm{L^{-1}f-G_\tau^{\mathrm{MTT}}f}
        \le
        \norm{L^{-1}}\,
        \bigl(1-e^{-\tau\Lambda_{\mathrm{coh}}}\bigr)\norm{f}.
\]
\end{proposition}

\begin{proof}
The first estimate follows from
\[
        L^{-1}-L^{-1}B_\tau=L^{-1}(I-B_\tau).
\]
The second estimate uses \cref{thm:sector}.
\end{proof}

\subsection{Canonical commutators}

The sharp canonical commutator
\[
        [\phi(x),\pi(y)]=i\hbar\delta(x-y)
\]
is replaced inside the coherent chart by
\[
        [\phi_{\mathrm{coh}}(x),\pi_{\mathrm{coh}}(y)]
        =
        i\hbar K_{\mathrm{MTT}}(x,y;\tau).
\]
The correction is controlled by the same sector error estimates. This is not loss of locality
as arbitrary nonlocality; it is finite coherent-sector locality determined by \(A,P,\tau\).

\subsection{Contact vertices}

A local \(n\)-point contact vertex can be written schematically as
\[
        V_\delta^{(n)}(x_1,\ldots,x_n)
        =
        \int_X\prod_{j=1}^n \delta(x-x_j)\,\dd x.
\]
The fixed-point replacement is
\[
        V_{\mathrm{MTT}}^{(n)}(x_1,\ldots,x_n;\tau)
        =
        \int_X\prod_{j=1}^n K_{\mathrm{MTT}}(x,x_j;\tau)\,\dd x.
\]
Thus the contact interaction is not softened by an arbitrary profile; it is softened by the
coherent kernel selected by the same fixed-point data.

\subsection{Measurement effects}

The distributional position effect
\[
        |x\rangle\langle x|
\]
is replaced by a finite coherent effect with kernel
\[
        E_x^{\mathrm{MTT}}(y,z;\tau)
        =
        K_{\mathrm{MTT}}(y,x;\tau)K_{\mathrm{MTT}}^\ast(z,x;\tau).
\]
The width of the measurement effect is therefore controlled by the proper-time/coherence scale
and spectral structure of the fixed-point sector.

\subsection{Noise kernels}

A temporal fixed-point operator \(A_t\) with memory scale \(\tau\) selects the correlation
kernel
\[
        C_{\mathrm{MTT}}(t,s)
        =
        \langle t|e^{-\tau A_t}|s\rangle
\]
rather than an exact white-noise delta. In Markov limits this may approach
\[
        D\delta(t-s),
\]
but the finite-memory kernel is the execution-level object.

\section{Worked example: circle with massive operator}

Let \(X=S^1\) with coordinate \(\theta\in[0,2\pi)\), and take
\[
        A_m=-\partial_\theta^2+m^2,\qquad m>0.
\]
The eigenfunctions are
\[
        \phi_n(\theta)=\frac{1}{\sqrt{2\pi}}e^{in\theta},
        \qquad
        \lambda_n=n^2+m^2.
\]
Let
\[
        P_N=\sum_{|n|\le N}|\phi_n\rangle\langle\phi_n|.
\]
Then
\[
        K_{N,\tau}^{\mathrm{MTT}}(\theta,\theta')
        =
        \frac{1}{2\pi}
        \sum_{|n|\le N}
        e^{-\tau(n^2+m^2)}e^{in(\theta-\theta')}.
\]
For fixed \(N\), this is a smooth finite kernel. For \(N\to\infty\), it becomes the heat
kernel of \(A_m\):
\[
        K_{\infty,\tau}^{\mathrm{MTT}}(\theta,\theta')
        =
        e^{-\tau m^2}
        \frac{1}{2\pi}
        \sum_{n\in\bbZ}e^{-\tau n^2}e^{in(\theta-\theta')}.
\]
As \(\tau\downarrow0\), this tends distributionally to the periodic delta
\[
        \delta_{S^1}(\theta-\theta').
\]

The finite correction to a coherent-band identity is explicit:
\[
        \norm{P_Nf-B_{N,\tau}f}_{L^2}^2
        =
        \sum_{|n|\le N}
        \left(1-e^{-\tau(n^2+m^2)}\right)^2|f_n|^2.
\]
Thus the correction is mode-resolved and directly computable from the fixed-point spectral
data.

\section{Programmatic consequence}

The previous delta papers established the forward structural dictionary:
\[
        \text{finite kernel}\longrightarrow \delta.
\]
This paper supplies the fixed-point reverse direction:
\[
        (A,P,\tau,\lambda^\ast)
        \longrightarrow
        K_{\mathrm{MTT}}
        \longrightarrow
        \text{finite correction}.
\]

\begin{quote}
\textbf{Execution principle.}
In any admissible fixed-point chart, replace delta distributions not by arbitrary mollifiers,
but by the projected heat/proper-time kernel selected by the fixed-point operator and
coherent projector.
\end{quote}

This is the point at which the delta dictionary becomes predictive in principle. Once a
sector supplies numerical or geometric data for \(A\), \(P\), \(\tau\), and \(\lambda^\ast\), the
corrections to point sources, commutators, contact interactions, measurement effects, and
noise kernels are fixed.

\section{Conclusion}

The central problem after the delta--projection dictionary was not whether finite kernels
can converge to Dirac deltas. They can. The execution problem was whether MTT selects a
specific finite kernel.

In the fixed-point regime, the answer is yes. The fixed-point data determine
\[
        B_\tau=P e^{-\tau A}P
\]
and hence
\[
        K_{\mathrm{MTT}}(x,y;\tau)
        =
        \langle x|Pe^{-\tau A}P|y\rangle.
\]
This kernel is smooth for positive proper time, covariant under admissible re-encoding,
identity-like on the coherent sector with explicit error bounds, suppresses incoherent modes
by the spectral gap, and tends to the Dirac delta only under the idealizing limit of complete
bandwidth and vanishing proper-time width.

Thus the program is no longer merely:
\[
        \delta=\text{singular limit of some finite object}.
\]
It becomes:
\[
        \boxed{
        \text{MTT fixed-point data select the finite object.}
        }
\]

\begin{thebibliography}{9}

\bibitem{Kato}
T. Kato,
\emph{Perturbation Theory for Linear Operators},
Springer, 1995.

\bibitem{ReedSimon}
M. Reed and B. Simon,
\emph{Methods of Modern Mathematical Physics I: Functional Analysis},
Academic Press, 1980.

\bibitem{Davies}
E. B. Davies,
\emph{Heat Kernels and Spectral Theory},
Cambridge University Press, 1989.

\bibitem{Taylor}
M. E. Taylor,
\emph{Partial Differential Equations II: Qualitative Studies of Linear Equations},
Springer, 2011.

\bibitem{HillePhillips}
E. Hille and R. S. Phillips,
\emph{Functional Analysis and Semi-Groups},
American Mathematical Society, 1957.

\end{thebibliography}

\end{document}
